\section{Introduction}\label{sec:intro}

Genotoxicity is a critical toxicological endpoint that assesses the potential of chemicals to induce DNA damage, 
mutations, or carcinogenesis and plays a central role in regulatory chemical evaluation \cite{Karamertzanis2025, Hayashi2022, EFSA2018}.
Within the OECD genetic toxicology framework, the bacterial reverse mutation assay (Ames test; OECD Test Guideline 471) 
continues to serve as a cornerstone for detecting point mutations, while mammalian \textit{in vitro} and \textit{in vivo} 
assays provide complementary evidence for broader forms of genetic damage. In parallel, ICH S2(R1) recommends that bacterial, 
mammalian \textit{in vitro}, and follow-up \textit{in vivo} studies should be interpreted within an integrated testing 
strategy rather than treated as interchangeable endpoints. Accordingly, genotoxicity prediction is not a 
single-endpoint problem, and differences in assay biology, protocol design, and regulatory interpretation must be considered 
when constructing computational datasets and models \cite{OECDOverview2017,OECD4712020,ICHS2R1}.

%규제 체계에서 시험을 의무화함에 따라 유전독성 데이터의 활용 가능성이 많아졌고, 이는 In silico model 개발을 촉진했다\citep{HASSELGREN2019104403}
The rapid expansion of publicly accessible toxicological data and the widespread adoption of machine learning 
have stimulated intense interest in \textit{in silico} genotoxicity prediction\cite{HASSELGREN2019104403}.

However, in this area, predictive performance is often constrained less by algorithmic sophistication than by the quality, 
provenance, and consistency of the underlying data. Classical cheminformatics studies have repeatedly shown that 
unresolved salts and mixtures, inconsistent parent--child representations, duplicate structures, erroneous structure records, 
and conflicting annotations can materially distort descriptor calculation, training distributions, and external validation 
results. Subsequent work has extended this argument from chemical structure curation to integrated chemical--biological 
curation, emphasizing that reproducible modeling requires explicit workflows for detecting duplicate records, 
reconciling inconsistent labels, and standardizing structure representations before any model development begins \cite{Fourches2010,Fourches2016}. 
For assay-driven toxicology datasets, the problem is even more acute: endpoint definitions, activity thresholds, 
cytotoxicity filters, signal interference, and purity criteria can all alter what is ultimately learned by a statistical model \cite{Nikolov2023}.

These issues are particularly relevant in genotoxicity, where historical databases frequently aggregate results generated 
across different laboratories, time periods, assay variants, and interpretive frameworks. Recent analyses of guideline-conform 
genotoxicity data have shown that substantial variability and uncertainty remain even within reference datasets assembled 
from standardized test methods, implying that the apparent performance ceiling of new approach methodologies 
should be interpreted in the context of intrinsic experimental noise rather than assumed ground truth. More broadly, 
recent discussions of QSAR predictivity in toxicology have highlighted recurring sources of poor model performance, 
including deficiencies in data quality, endpoint consistency, mechanistic relevance, uncertainty characterization, 
and applicability-domain assessment \cite{Raitano2026,Cronin2025}. In other words, improving genotoxicity prediction 
requires more than a better classifier; it requires a transparent and endpoint-aware data policy.

%현재 범용적인 QSAR 모델인 OECD QSAR Toolbox와 VEGA-QSAR는 산업용 화학물질에 적용 도메인이 알맞지 않음\citep{YOON2026100405} 추가할까 고민 중

This need for transparency aligns closely with the FAIR principles, which call for scientific data and associated workflows 
to be \textbf{F}indable, \textbf{A}ccessible, \textbf{I}nteroperable, and \textbf{R}eusable \cite{Wilkinson2016}. 
In toxicology, FAIR implementation has been framed not only as a data-management ideal but also as a prerequisite for 
integrating legacy hazard data, new approach methodologies, and computational models into coherent evidence streams \cite{Watford2019}. 
Recent work has further argued that \textit{in silico} toxicology models themselves should be made FAIR, 
with explicit documentation of endpoint definition, model provenance, reporting standards, and conditions of reuse to 
improve reproducibility and regulatory uptake \cite{Cronin2023}. This perspective is consistent with 
the OECD (Q)SAR Assessment Framework, which emphasizes correct model input, applicability domain, prediction reliability,
and fitness for purpose, and with recent genotoxicity-specific guidance that stresses the importance of defined endpoints, 
well-documented predictions, and domain-aware interpretation when applying (Q)SAR approaches in regulatory contexts \cite{OECDQAF2024,COM2026}.

Against this background, data curation and preprocessing are not merely technical preliminaries but scientific decisions 
that shape the scope and validity of downstream inference. Ames mutagenicity is comparatively well supported by public benchmark
resources, which has facilitated broad model development and external comparison \cite{Hansen2009}. 
By contrast, mammalian \textit{in vitro} and \textit{in vivo} genotoxicity datasets are often assembled from more 
heterogeneous evidence streams and therefore place greater demands on explicit endpoint definition, provenance tracking, 
and rule-based curation \cite{Nikolov2023,Raitano2026,COM2026}. This distinction motivates an endpoint-specific approach 
in which structure standardization, salt and metal handling, duplicate resolution, label harmonization, and source 
annotation are treated as study variables rather than hidden preprocessing defaults.

In the present study, we started from an integrated source workbook that preserved both original test outcomes and identifying metadata, and we used this resource to construct endpoint-specific curated datasets for genotoxicity modeling. For Ames mutagenicity, the final curated dataset combined records originating from the raw industrial corpus with benchmark-derived drug records, while explicitly retaining provenance through a domain variable. In contrast, the mammalian \textit{in vitro} and \textit{in vivo} datasets were curated only from the original source without external benchmark augmentation. We therefore aimed not merely to build another predictive model, but to examine how transparent curation and preprocessing choices affect dataset integrity, comparability across endpoints, and the reliability of downstream \textit{in silico} analyses. Framed in this way, the study addresses a practical question that is highly relevant to both cheminformatics and regulatory toxicology: whether genotoxicity datasets become more reproducible, interpretable, and FAIR-aligned when curation rules are made endpoint-specific, provenance-aware, and explicitly reportable.